\section{The adjoint-addition principle}\label{addition:section}

The section comparison reduces subadditivity to a statement about
$K_W+T+M$. We prove that statement here from the two positivity
inputs specified below. Its two maps have different purposes:
$h$ is the fibration whose fiber dimension is to be added, whereas
$p$ expresses the nef term as a pullback from a projective base.

The first input is the precise adjoint-positivity theorem for the
line produced by the construction. We use the highest-step and
parabolic-extension conventions of Section~\ref{rankone:section}.

\begin{theorem}[Adjoint-positivity input]\label{period:adjoint}
Let $B$ be a smooth compact connected manifold in class $\mathcal C$
and $B^\circ\subset B$ a dense Zariski open set with SNC complement.
Let $\mathbb V$ be a pure real-polarizable variation of Hodge
structures on $B^\circ$ with an integral ambient lattice. Suppose
a complex direct summand of $\mathbb V_{\mathbf C}$ has a
rank-one highest nonzero Hodge step, with parabolic rational
extension $L$. There are a modification $\tau:B'\to B$, a
surjective morphism $p:B'\to S$ with connected fibers to a smooth
projective variety, and a nef rational line bundle $P$ on $S$ such that
\begin{equation}\label{period:conclusion}
 \tau^*L=p^*P,\qquad K_S+aP\text{ is big for some }a\in\mathbf Q_{>0}.
\end{equation}
Positive rational rescaling of $L$ is allowed, and the line identity
persists under further modifications. If $S$ is a point, $\tau^*L$
is rationally trivial.
\end{theorem}

The full proof is in Section~\ref{period:section}. The lattice belongs to the ambient
variation: the complex summand need not itself be defined over
$\mathbf Q$, and the real polarization need not be rational.

The second input supplies addition when the fiber already has a big
log canonical divisor. It will be proved in Section~\ref{generaltype:section}.

\begin{proposition}\label{generaltype:addition}
Let \(h\colon V\to Y\) be a surjective holomorphic map with connected fibers
between smooth compact complex manifolds in Fujiki class \(\mathcal C\).
Let \(T\) and \(C\) be effective rational divisors with simple normal crossing
support and coefficients in \([0,1]\) on \(V\) and \(Y\), respectively.
Suppose that
\begin{enumerate}[label=\textup{(\roman*)}]
\item \(K_F+T_F\) is big for a very general smooth fiber \(F\) of \(h\);
\item \(C\leq B(h,T)\);
\item \(T_E=1\) for every prime divisor \(E\) of \(V\) whose image in \(Y\)
has codimension at least two.
\end{enumerate}
Then
\begin{equation}\label{generaltype:inequality}
  \kappa(V,K_V+T)\ \geq\
  \dim F+\kappa(Y,K_Y+C).
\end{equation}
As usual, the assertion is automatic if the last term is \(-\infty\).
\end{proposition}


General-type addition will give sections after a positive ample
twist from $S$. To remove that twist, we will need one fixed section
with a negative ample twist. The next lemma produces it from a big
base divisor, once nonvanishing on the general $p$-fiber is known.
It uses weak positivity with exceptional poles and
descent to a fixed smooth reference, following the method of
\cite[Remark~6.5(3) and Proposition~6.6]{CampanaOrbifolds}.
We include the descent step explicitly.

\begin{lemma}[Effectivity after a big base twist]
\label{application:weak-effectivity}
Let $p:(W,T)\to S$ be surjective, with $W$ smooth compact in
class $\mathcal C$, $S$ smooth projective, and $T$ a rational SNC
boundary. Suppose one fixed divisible relative log pluricanonical
system is nonzero on very general fibers. For every big rational
line bundle $H$ on $S$, the rational line bundle
$K_{W/S}+T+p^*H$ has a nonzero section in some positive degree.
\end{lemma}

\begin{proof}
Choose a K\"ahler source modification, flatten over a smooth
projective modification $\pi:S'\to S$, and resolve the flat main
transform. Write $\mu:W'\to W$ and $p':W'\to S'$ for the
resulting maps. Give $W'$ the strict transform of $T$ plus the
reduced $\mu$-exceptional divisor, denoted $T'$. Then
\begin{equation}\label{application:source-discrepancy}
 K_{W'}+T'=\mu^*(K_W+T)+E,
 \qquad E\geq0\ \text{and $\mu$-exceptional}.
\end{equation}
Every source prime over a subset of codimension at least two in
$S'$ is exceptional over the fixed $W$: before resolving the
source, flatness bounds that inverse image by codimension at least
two. Put $R=K_{S'/S}\geq0$ and $H'=\pi^*H+R$. The exact identity
\begin{equation}\label{application:twist-cancellation}
 K_{W'/S'}+T'+p'^*H'
 =\mu^*(K_{W/S}+T+p^*H)+E
\end{equation}
will allow descent.

Choose a divisible $m>0$ such that
$\mathcal E=p'_*\mathcal O_{W'}(m(K_{W'/S'}+T'))$ has positive
rank. Generic base change and the hypothesis provide such an $m$.
Fujino's weak positivity theorem applies to this log canonical pair
and its projective base \cite[Theorem~1.1]{FujinoWeakPositivity}.
Choose an ample Cartier divisor $A_0$ and an integer $\alpha>0$
such that $H'-A_0/(m\alpha)$ is big. For some integer $\beta>0$,
\[
 (\operatorname{Sym}^{\alpha\beta}\mathcal E)^{**}
        \otimes\mathcal O_{S'}(\beta A_0)
\]
is generated by global sections at the generic point.

On the locally free locus of $\mathcal E$, multiplication of
relative sections gives an evaluation map to
\[
 \mathcal O_{W'}\bigl(m\alpha\beta(K_{W'/S'}+T')
                         +\beta p'^*A_0\bigr).
\]
It is nonzero generically, because a nonzero relative section has
nonzero powers. Generic generation supplies a global section whose
image is nonzero. The complement of the locally free locus has
codimension at least two. Multiply a sufficiently divisible power
of the evaluated section by a section of a sufficiently divisible
multiple of
$m\alpha\beta\,p'^*(H'-A_0/(m\alpha))$.
The latter exists by bigness. This produces a nonzero section of a
multiple of the left side of \eqref{application:twist-cancellation}, except possibly
at source primes above that codimension-two set.

Locally, a meromorphic frame of $\mathcal E$ expresses the evaluated
section as a meromorphic combination of products of relative
sections. Thus the resulting form is meromorphic, and its only
possible divisorial poles lie above that omitted codimension-two set.
Those primes, and $E$, are exceptional over $W$. At every prime
of the original smooth $W$, the resulting meromorphic form therefore
has the required order for a multiple of $K_{W/S}+T+p^*H$.
Divisorial regularity and extension across codimension two give a
global section on $W$. If $S$ is a point, the same conclusion is
the assumed nonvanishing.
\end{proof}

\begin{proposition}[Adjoint addition]
\label{addition:principle}
Let $h:W\to Y$ be a fibration between smooth compact manifolds in
class $\mathcal C$. Let $T,C$ be rational SNC boundaries with
$C\leq B(h,T)$, and suppose every prime of $W$ whose image has
codimension at least two in $Y$ has coefficient one in $T$.
Let $p:W\to S$ be a surjective morphism with connected fibers to
a smooth projective variety. Suppose $P$ is a nef rational line
bundle on $S$, $M=p^*P$, and $K_S+a_0P$ is big for some rational
$a_0>0$. If $K_W+T+M$ is big on very general $h$-fibers, then
\[
 \kappa(W,K_W+T+M)
 \geq\dim(W/Y)+\kappa(Y,K_Y+C).
\]
\end{proposition}

\begin{proof}
The assertion is automatic if the last term is $-\infty$.
Assume it is nonnegative, put $D=K_W+T$, and write
$k=\dim(W/Y)$. We may first replace $W$ by a K\"ahler
modification, adding the reduced exceptional divisor to its strict
boundary. The log section spaces, also after tensoring by the
pullback of $M$, are unchanged. Fiberwise bigness and the
inf-multiplicity boundary are preserved; newly exceptional primes
have coefficient one. Thus all hypotheses remain valid.

If $S$ is a point, $M$ is rationally trivial and
\Cref{generaltype:addition} proves the proposition. Suppose $\dim S>0$,
and choose an ample Cartier divisor $A$ on $S$. A very general
$h$-fiber is K\"ahler and carries the big rational line
$(D+M)|_{W_y}$, so it is projective. Openness of its big cone
gives a rational $c$ with $0<c<1$ for which
$(D+cM)|_{W_y}$ is still big. This choice works on very general
fibers: choose the first fiber outside the countably many
cohomology and evaluation-rank exceptional loci for rational $c$
and divisible degrees. The one full-rank system just found then
has the same rank generically. For zero-dimensional fibers this
argument is unnecessary.

For every positive rational $\eta$, $cP+\eta A$ is ample.
Choose a sufficiently divisible large $N$ and a general member
of the basepoint-free system $|Np^*(cP+\eta A)|$. Dividing by
$N$ gives an effective rational divisor $Q_\eta$ such that
\[
 Q_\eta\sim_{\mathbf Q}cM+\eta p^*A,
 \qquad T+Q_\eta\ \text{is an SNC boundary}.
\]
Bertini is applied simultaneously to the finitely many strata of
$T$; the member contains no old boundary component, and increasing
$N$ bounds its new coefficients by one. Existing coefficient-one
components remain unchanged, and increasing the source boundary
can only increase $B(h,T)$. Thus all boundary hypotheses of
\Cref{generaltype:addition} hold for $T+Q_\eta$. Its fiber adjoint is big
by the choice of $c$ and nefness of the added ample pullback.
Consequently
\begin{equation}\label{application:perturbed-bound}
 \kappa(W,D+cM+\eta p^*A)
       \geq k+\kappa(Y,K_Y+C)
       \qquad(\eta\in\mathbf Q_{>0}).
\end{equation}

The systems in \eqref{application:perturbed-bound} already have the required
image dimension. It remains to remove their ample term without
losing these ratios of sections. In particular one such divisor has
a nonzero section in a fixed
positive degree. Restriction of that section to a very general
$p$-fiber is nonzero; both pulled-back twists become trivial
there. Hence
\begin{equation}\label{application:p-fiber-nonvanishing}
 \kappa(W_s,(K_W+T)|_{W_s})\geq0
 \quad\text{for very general }s\in S.
\end{equation}
One fixed relative log pluricanonical direct image for
$p:(W,T)\to S$ therefore has positive rank.

Choose rational $a>\max\{1,a_0\}$. Since $P$ is nef,
$K_S+aP$ is big. For sufficiently small rational $\lambda>0$,
$H=K_S+aP-\lambda A$ is still big.
Lemma~\ref{application:weak-effectivity} now gives a nonzero section of a
positive multiple of
\begin{equation}\label{application:negative-twist}
 E_-=D+aM-\lambda p^*A
     =K_{W/S}+T+p^*H.
\end{equation}
Put
\[
 \theta=\frac{1-c}{a-c},
 \qquad \eta=\frac{(1-c)\lambda}{a-1},
 \qquad E_+=D+cM+\eta p^*A.
\]
Then $0<\theta<1$, $\eta>0$, and direct calculation gives
\begin{equation}\label{addition:cancel}
 D+M=(1-\theta)E_++\theta E_-.
\end{equation}
Fix a nonzero section $e$ of $rE_-$ for one denominator-clearing
$r>0$. For sufficiently divisible $n$, multiplication by
$e^{n\theta/r}$ maps
\[
 H^0(W,n(1-\theta)E_+)
       \lhook\joinrel\longrightarrow H^0(W,n(D+M)).
\]
The exponent is an integer, and the source degrees range through a
Veronese subsequence. On the complement of the zero divisor of $e$,
every ratio of two sections is unchanged. These systems therefore
have the same image dimensions before and after multiplication.
Equation~\eqref{application:perturbed-bound} for the displayed value of $\eta$
proves the asserted lower bound.
\end{proof}
